% ------------------------------------------------------------------------------
% Metodologia
% ------------------------------------------------------------------------------

\chapter{Metodologia}
\label{chap_metodologia}

Este capítulo descreve os procedimentos metodológicos adotados.
Trata-se de um estudo quantitativo, observacional e longitudinal, baseado exclusivamente em dados públicos.
A \autoref{sec_fontes_dados} apresenta as fontes de dados; a \autoref{sec_coleta_e_tratamento_de_dados} descreve o \textit{pipeline} de coleta e tratamento; a \autoref{sec_desidentificacao} discute a desidentificação dos dados e as variáveis \textit{proxy}; e a \autoref{sec_analise_estatistica} detalha a especificação econométrica e os procedimentos de inferência.

\section{Fontes de dados}
\label{sec_fontes_dados}

O estudo integra quatro fontes oficiais, descritas no \autoref{qua_fontes}.
Todos os arquivos foram obtidos nos repositórios de dados abertos dos respectivos órgãos e armazenados em sua forma original antes de qualquer transformação.

\begin{quadro}[!htb]
    \centering
    \caption{Fontes de dados utilizadas}
    \label{qua_fontes}
    \footnotesize
    \begin{tabular}{|p{3.0cm}|p{4.6cm}|p{2.2cm}|p{4.2cm}|}
        \hline
        \textbf{Fonte} & \textbf{Conteúdo} & \textbf{Período} & \textbf{Variáveis extraídas} \\
        \hline
        CGU -- Portal da Transparência & Folhas mensais de pagamento do Pé-de-Meia (18 colunas, codificação \textit{latin-1}) & Mar./2024 a mar./2025 & Valor e quantidade de parcelas por UF e tipo de incentivo \\
        \hline
        INEP -- Indicadores Educacionais & Taxas de rendimento escolar por UF e dependência administrativa & 2022 a 2025 & Taxas de aprovação, reprovação e abandono do ensino médio \\
        \hline
        INEP -- IDEB/SAEB & Planilhas de divulgação do IDEB do ensino médio por UF e rede & 2019, 2021, 2023 e 2025 & Notas do SAEB, $N$, $P$ e IDEB observado \\
        \hline
        CadÚnico (Portal de Dados Abertos de Alagoas) & Estoque mensal de famílias e pessoas cadastradas por UF & Dez./2024 & Famílias e pessoas cadastradas \\
        \hline
    \end{tabular}
    \fonte{Elaborado pelo autor (2026) com base em \citeonline{CGU2026PeDeMeia}, \citeonline{INEP2026Rendimento}, \citeonline{INEP2026IDEB} e \citeonline{Alagoas2026CadUnico}.}
\end{quadro}

\section{Coleta e tratamento de dados}
\label{sec_coleta_e_tratamento_de_dados}

O \textit{pipeline} foi implementado em Python, com as bibliotecas \textit{pandas} \cite{McKinney2010} para manipulação de dados, \textit{statsmodels} \cite{Seabold2010} para a estimação econométrica e \textit{Matplotlib} \cite{Hunter2007} e \textit{seaborn} \cite{Waskom2021} para as visualizações.
O fluxo de processamento, ilustrado na \autoref{fig_pipeline}, divide-se em três etapas: coleta automatizada dos arquivos brutos, tratamento e consolidação das bases e, por fim, modelagem e visualização.

\begin{figure}[!htb]
    \centering
    \caption{\textit{Pipeline} de coleta, tratamento e análise dos dados}
    \label{fig_pipeline}
    \begin{tikzpicture}[
        node distance=0.55cm and 0.9cm,
        caixa/.style={draw, rounded corners, align=center, font=\footnotesize,
                      minimum width=3.6cm, minimum height=1.0cm, fill=gray!8},
        seta/.style={-{Stealth[length=2.2mm]}, thick}
    ]
        \node[caixa] (cgu)  {CGU\\Folhas mensais};
        \node[caixa, below=of cgu]  (rend) {INEP\\Taxas de rendimento};
        \node[caixa, below=of rend] (ideb) {INEP\\IDEB/SAEB};
        \node[caixa, below=of ideb] (cad)  {CadÚnico\\Famílias e pessoas};

        \node[caixa, right=1.3cm of $(rend)!0.5!(ideb)$, minimum height=2.2cm] (mat)
            {Matriz mestre\\324 linhas $\times$ 39 colunas\\(27 UFs $\times$ 2 redes $\times$ 6 anos)};

        \node[caixa, right=1.3cm of mat, yshift=0.9cm]  (did) {Modelos DiD\\(21 regressões)};
        \node[caixa, right=1.3cm of mat, yshift=-0.9cm] (graf) {Tabelas e gráficos};

        \foreach \n in {cgu, rend, ideb, cad} \draw[seta] (\n.east) -- (mat.west);
        \draw[seta] (mat.east) -- (did.west);
        \draw[seta] (mat.east) -- (graf.west);
    \end{tikzpicture}
    \fonte{Elaborada pelo autor (2026).}
\end{figure}

Na etapa de tratamento, os pacotes mensais da CGU foram lidos integralmente e agregados por UF, mês de folha e tipo de incentivo.
Foram processadas todas as nove competências de 2024 (março a dezembro, inexistindo a folha de julho em razão do recesso escolar), além das folhas de janeiro e fevereiro de 2025, que contêm os depósitos de conclusão e de Enem referentes ao ano letivo de 2024.
As planilhas do INEP foram filtradas para o ensino médio, para o nível de agregação UF e para as dependências administrativas estadual e privada.

As bases foram então integradas em uma matriz mestre balanceada de 324 observações (27 UFs $\times$ 2 redes $\times$ 6 anos: 2019, 2021, 2022, 2023, 2024 e 2025) e 39 colunas, documentada no \autoref{chap_apendice_dicionario}.
Os valores ausentes são estruturais: os indicadores do SAEB e do IDEB existem apenas em anos ímpares, as taxas anuais de rendimento foram consolidadas para 2022 a 2025 e as variáveis financeiras do Pé-de-Meia referem-se exclusivamente às 27 observações da rede estadual em 2024.

\subsection{Auditoria de consistência das folhas de pagamento}
\label{subsec_auditoria}

Durante a coleta, identificou-se que a versão do arquivo de março de 2024 inicialmente obtida continha 4.205.138 parcelas, totalizando R\$ 841.027.600,00.
A versão retificada posteriormente publicada pela CGU contém 2.416.252 parcelas e R\$ 483.250.400,00, uma redução de 1.788.886 registros ($-42{,}54\%$).
A inspeção dos beneficiários maiores de 18 anos, para os quais o NIS é divulgado, mostrou que 68,55\% dos identificadores apresentavam duas linhas idênticas na mesma data de pagamento na versão antiga, o que caracteriza duplicação.
Todas as análises deste trabalho utilizam exclusivamente a versão retificada.

\section{Desidentificação dos dados e variáveis \textit{proxy}}
\label{sec_desidentificacao}

Na folha retificada de março de 2024, 88,86\% dos registros trazem o nome do beneficiário substituído pela expressão ``beneficiário menor de 18 anos'', com CPF e NIS ausentes.
Essa desidentificação decorre das normas de proteção de dados de crianças e adolescentes previstas na LGPD \cite{Brasil2018LGPD} e no ECA \cite{Brasil1990ECA}.
Como consequência, não há chave individual confiável para rastrear um mesmo estudante ao longo das folhas mensais.

Por esse motivo, adotaram-se duas variáveis \textit{proxy}.
Como o Incentivo-Matrícula é pago uma única vez por estudante elegível em cada ano letivo, a contagem de parcelas dessa modalidade em 2024 (4.008.269) foi utilizada como estimativa do número de estudantes beneficiários.
Analogamente, a contagem de parcelas de Incentivo-Conclusão referentes ao ano letivo de 2024 pagas até fevereiro de 2025 (3.037.316) foi utilizada como estimativa do número de estudantes que receberam o depósito de conclusão até essa data.

\section{Análise estatística}
\label{sec_analise_estatistica}

\subsection{Especificação do modelo de Diferenças em Diferenças}
\label{subsec_especificacao}

O grupo tratado é a rede pública estadual ($D_r = 1$), e o grupo de comparação é a rede privada ($D_r = 0$), não elegível ao programa.
O período pós-programa é definido por $\mathit{Post}_t = 1$ para $t \geq 2024$.
A especificação principal é dada pela \autoref{eq_did}:
\begin{equation}
    \label{eq_did}
    Y_{urt} = \alpha_{ur} + \lambda_t + \beta\,(D_r \times \mathit{Post}_t) + \varepsilon_{urt},
\end{equation}
em que $Y_{urt}$ é o desfecho da UF $u$, na rede $r$ e no ano $t$; $\alpha_{ur}$ é o efeito fixo do par UF $\times$ rede, que absorve todas as características invariantes no tempo de cada combinação; $\lambda_t$ é o efeito fixo de ano, que absorve choques comuns às duas redes; e $\beta$ é o coeficiente de interesse, que mede a variação diferencial da rede estadual no período pós-programa.

Os desfechos e as janelas analíticas são os seguintes:
\begin{itemize}
    \item taxas anuais de abandono, aprovação e reprovação do ensino médio, no painel 2022--2025 ($n = 216$);
    \item notas do SAEB em Matemática e Língua Portuguesa, nota padronizada $N$, indicador de rendimento $P$ e IDEB observado, no painel bienal completo 2019--2025 ($n = 216$) e na janela imediata 2023 \textit{versus} 2025 ($n = 108$).
\end{itemize}

\subsection{Testes de pré-tendência}
\label{subsec_pretendencia}

Para avaliar a plausibilidade das tendências paralelas, estimou-se, restrito ao período pré-programa, o modelo da \autoref{eq_pretendencia}:
\begin{equation}
    \label{eq_pretendencia}
    Y_{urt} = \alpha_{ur} + \lambda_t + \sum_{k \in \mathcal{K}} \theta_k\,\big(D_r \times \mathbf{1}[t = k]\big) + \varepsilon_{urt},
\end{equation}
em que $\mathcal{K} = \{2023\}$ para as taxas de rendimento (referência: 2022) e $\mathcal{K} = \{2021, 2023\}$ para os indicadores do SAEB e do IDEB (referência: 2019).
Sob tendências paralelas, espera-se $\theta_k = 0$ para todo $k$.

\subsection{Inferência}
\label{subsec_inferencia}

Os modelos foram estimados por Mínimos Quadrados Ordinários, com erros-padrão robustos agrupados por UF ($G = 27$ \textit{clusters}), o que acomoda a correlação serial e a correlação entre as duas redes de uma mesma UF \cite{Bertrand2004}.
Dado o número reduzido de \textit{clusters}, os testes de hipótese utilizam a distribuição $t$ de Student com $G - 1 = 26$ graus de liberdade \cite{Cameron2015}.

No total, foram estimados 21 modelos, que produzem 26 testes de hipótese sobre coeficientes de interação: 13 testes DiD pós-programa (3 de fluxo anual, 5 no painel bienal completo e 5 na janela imediata) e 13 testes de pré-tendência.
Para controlar a taxa de erro do tipo I decorrente da multiplicidade de testes, os resultados são também confrontados com o limiar de Bonferroni \cite{Dunn1961}, definido pela \autoref{eq_bonferroni}:
\begin{equation}
    \label{eq_bonferroni}
    \alpha_{\text{Bonf}} = \frac{\alpha}{m},
\end{equation}
em que $\alpha = 0{,}05$ e $m$ é o número de testes da família considerada.
Para os 13 testes DiD, $\alpha_{\text{Bonf}} \approx 0{,}0038$; para os 5 desfechos da janela imediata, $\alpha_{\text{Bonf}} = 0{,}0100$.
Como análise de robustez, os modelos foram também estimados em uma especificação preliminar sem efeitos fixos de unidade e com distribuição normal, cujos resultados constam no \autoref{chap_apendice_comparacao}.
